\documentclass[10pt,a4paper]{amsart}
\usepackage[margin=19mm]{geometry}
\usepackage{amsmath,amssymb,mathtools}
\usepackage[scr=rsfs,cal=euler]{mathalfa}
\usepackage{xfrac}
\usepackage[unicode,hidelinks,bookmarks=false]{hyperref}
\newcommand{\ie}{i.e.\ }
\newcommand{\cf}{cf.\ }
\newcommand{\resp}{respectively\ }
\newcommand{\Subsection}{\subsection}
\providecommand{\nobreakdash}{\nobreak\hbox{-}}
\setlength{\emergencystretch}{2em}
\multlinegap0pt
\usepackage{mathtools}
\usepackage{xcolor}
\usepackage[normalem]{ulem}
\usepackage{amssymb}
\usepackage{braket}
\usepackage[shortlabels]{enumitem}
\newtheorem{lemma}{Lemma}
\newtheorem{proposition}{Proposition}
\def\CA{{\mathcal A}}
\def\CH {{\mathcal H}}
\def\La{{\Lambda}}
\def\hl{{\hat \lambda}}
\def\l{{\lambda}}
\def\lab{\operatorname{lab}}
\def\tr{\operatorname{tr}}
\title{A Classification of Invertible Stabilizer Codes}
\author{Roman Geiko, Georgii Shuklin}
\date{Source: arXiv:2512.01142v2 (CC BY 4.0)}
\begin{document}
\maketitle

\noindent\emph{Source and licence.} A Classification of Invertible Stabilizer Codes. Version arXiv:2512.01142v2 (CC BY 4.0), distributed by arXiv under the Creative Commons Attribution 4.0 International licence, http://creativecommons.org/licenses/by/4.0/, as recorded in the arXiv licence metadata for this exact version and shown on the abstract page https://arxiv.org/abs/2512.01142v2. Full licence notice: \texttt{licenses/arxiv-2512.01142-CC-BY-4.0.txt}.\par\smallskip

\noindent\textit{Editorial scope.} Counted unit: the article's proposition prop:recover\_code\_from\_hamiltonian (source0.tex lines 1125--1132). The article's complete original proof of it is delivered with it (source0.tex lines 1134--1237, one \texttt{proof} environment of 104 lines). No mathematics is written, repaired or re-derived: the statement and the proof are the article's own lines, in the article's own notation, and no retained line is substituted or re-flowed.  Retained uncounted as the source-local context the proof needs: lines 996--999; lines 1001--1010; lines 1068--1078.  Nothing was omitted from any retained span, so the package carries no omission record.  No retained line contains a citation, so the wrapper carries no bibliography and no bibliography entry.  No retained line contains a figure environment, an image-inclusion command or drawing code, so no image is delivered and the package declares no asset dependency.  Editorial formatting only, in addition: an AMS \texttt{amsart} preamble that restates the article's own declarations and macros used by the retained lines, namely the environments \texttt{lemma}, \texttt{proposition} on separate counters, together with the standard packages those lines need.  The retained environments print in this document as Lemma 1, Proposition 1 in order of appearance, whereas the article prints the same items with the numbering of its own full document.  The complete native source of the article as distributed in the arXiv e-print for this exact version is bundled unchanged as \texttt{sources/190104/source0.tex}.
\begin{align}\label{eq:stabilizer_hamiltonian}
    \CH_{\iota}
    =\sum_{j\in\La}\sum_{k=1}^{s}Q_{j,k}.
\end{align}

\begin{lemma}\label{lem:linear_independence_quasi_pauli}
Let $u_1,\ldots,u_N\in\textbf{P}(P,\l;L_1)$ have pairwise distinct
phase-free labels
\begin{align*}
    p_a\coloneqq\lab(u_a)\in P,
    \qquad 1\leqslant a\leqslant N.
\end{align*}
Then, $u_1,\ldots,u_N$ are linearly independent in
$\CA(P,\l;L_1)$.
\end{lemma}

\begin{proposition}\label{prop:derivation}
The formal sum in \eqref{eq:stabilizer_hamiltonian} defines a derivation
\begin{align}\label{eq:derivation}
    D_{\CH_{\iota}}(a)
    \coloneqq
    \sum_{j\in\La}\sum_{k=1}^{s}[Q_{j,k},a]
    \in\CA(P,\l;L_1),
    \qquad
    a\in\CA(P,\l;L_1).
\end{align}
\end{proposition}

\begin{proposition}\label{prop:recover_code_from_hamiltonian}
Let $\CH_{\iota}$ be the commuting-projector Hamiltonian constructed from a
stabilizer code $F\subset P$. Then
\begin{align}\label{eq:hamiltonian_centralizer_equals_orthogonal}
    K_{\CH_{\iota}}=F^{\perp}.
\end{align}
In particular, $K_{\CH_{\iota}}$ reproduces $F$ if $F$ is locally topologically ordered.
\end{proposition}

\begin{proof}
Let $u_p\in\textbf{P}(P,\l;L_1)$ have label $p\in P$, and define
\begin{align}
    \chi_p(g)
    \coloneqq
    e^{2\pi i\tr\{\hl(g,p)\}},
    \qquad g\in F.
\end{align}
Then
\begin{align}
    \iota(g)u_p=\chi_p(g)u_p\iota(g).
\end{align}
Writing
\begin{align}
    g_{j,k}\coloneqq x^jf_k,
\end{align}
the definition of $Q_{j,k}$ gives
\begin{align}\label{eq:commutator_projector_quasipauli}
    [Q_{j,k},u_p]
    =
    \frac{1}{n_k}
    \sum_{\ell=1}^{n_k-1}
    \bigl(1-\chi_p(\ell g_{j,k})\bigr)
    u_p\iota(\ell g_{j,k}).
\end{align}

For fixed $p$, the set
\begin{align}
    I_p
    \coloneqq
    \left\{
        (j,k)
        \,\middle|\,
        \chi_p(g_{j,k})\neq1
    \right\}
\end{align}
is finite by the same argument as in Proposition \ref{prop:derivation}.

Grouping the finitely many nonzero terms in
\eqref{eq:commutator_projector_quasipauli} according to their stabilizer
label gives
\begin{align}\label{eq:derivation_grouped_by_stabilizer_label}
    D_{\CH_{\iota}}(u_p)
    =
    \sum_{g\in F\setminus\{0\}}
    c_p(g)
    \bigl(1-\chi_p(g)\bigr)
    u_p\iota(g),
\end{align}
where only finitely many coefficients are nonzero and
\begin{align}
    c_p(g)
    \coloneqq
    \sum_{\substack{
        (j,k)\in I_p,\ 1\leqslant\ell<n_k\\
        \ell g_{j,k}=g
    }}
    \frac{1}{n_k}
    \geqslant0.
\end{align}

Suppose that $(j,k)\in I_p$ and set $g=g_{j,k}$. The term with
$\ell=1$ shows that
\begin{align}
    c_p(g)>0,
\end{align}
while
\begin{align}
    1-\chi_p(g)\neq0.
\end{align}
Moreover, as $g$ varies, the operators $u_p\iota(g)$ have the distinct
labels $p+g$ and are therefore linearly independent by Lemma \ref{lem:linear_independence_quasi_pauli}. Consequently,
\begin{align}
    D_{\CH_{\iota}}(u_p)=0
    \quad\Longleftrightarrow\quad
    \chi_p(x^jf_k)=1
\end{align}
for every $j\in\La$ and every $1\leqslant k\leqslant s$.

Equivalently,
\begin{align}
    \tr\{\hl(x^jf_k,p)\}=0
\end{align}
for every $j$ and $k$. Since the form is conjugate-linear in its first
variable,
\begin{align}
    \hl(x^jf_k,p)=x^{-j}\hl(f_k,p).
\end{align}
The traces of all monomial translates detect every coefficient of an element
of $S^{-1}R/R$. Therefore the preceding condition is equivalent to
\begin{align}
    \hl(f_k,p)=0
    \qquad
    \text{for every }k.
\end{align}
Since the elements $f_1,\ldots,f_s$ generate $F$ over $R$, this is
equivalent to
\begin{align}
    p\in F^{\perp}.
\end{align}
This proves
\eqref{eq:hamiltonian_centralizer_equals_orthogonal}. If $F$ is locally
topologically ordered, then $F^{\perp}=F$, and the final assertion follows.
\end{proof}

\end{document}
